\section{Πρόβλημα 2}

Θεωρούμε το εξής παίγνιο που μοντελοποιεί το δίλημμα του φυλακισμένου με τις
επιλογές $\{C, D\}$ για κάθε παίκτη:

\begin{center}
\begin{tabular}{ccc}
        & $C$  & $D$  \\
    $C$ & 2, 2 & 0, 3 \\
    $D$ & 3, 0 & 1, 1 \\
\end{tabular}
\end{center}

Ορίζουμε ένα νέο παίγνιο με βάση το παραπάνω, όπου η τελική ωφέλεια κάθε παίκτη
εξαρτάται και από την ωφέλεια που έχει ο άλλος παίκτης στο αρχικό παίγνιο.
Αν $(s, t)$ είναι το ζεύγος στρατηγικών που επιλέγουν οι 2 παίκτες, τότε η
τελική χρησιμότητα του παίκτη 1 είναι
$$
u'_1(s, t) = u_1(s, t) + \alpha u_2(s, t)
$$
όπου $\alpha$ μια παράμετρος με $\alpha \in [0, 1]$, και $u_1(s,t), u_2(s,t)$
είναι οι χρησιμότητες στο αρχικό παίγνιο. Ομοίως για τον παίκτη 2 έχουμε
$u'_2(s, t) = u_2(s, t) + \alpha u_1(s, t)$.

Το νέο παίγνιο ορίζεται ως εξής:
\begin{center}
\begin{tabular}{ccc}
        & $C$  & $D$  \\
    $C$ & $2 + 2 \alpha, 2\alpha + 2$ & $0 + 3\alpha, 0\alpha + 3$ \\
    $D$ & $3 + 0 \alpha, 3\alpha + 0$ & $1 + \alpha, \alpha + 1$ \\
\end{tabular}
\end{center}
Δηλαδή,
\begin{center}
\begin{tabular}{ccc}
        & $C$  & $D$  \\
    $C$ & $2 + 2 \alpha, 2\alpha + 2$ & $3\alpha, 3$ \\
    $D$ & $3, 3\alpha$ & $1 + \alpha, \alpha + 1$ \\
\end{tabular}
\end{center}


\subsection{Ερώτηση (i)}

Θα βρούμε για ποιες τιμές του $\alpha$ το νέο παίγνιο εξακολουθεί να αποτελεί
αναπαράσταση του διλήμματος του φυλακισμένου. Για να συνεχίζει το το παίγνιο να
αποτελεί αναπαράσταση του διλήμματος του φυλακισμένου, πρέπει να ισχύει ότι:
$$
u'_1(D, C) > u'_1(C, C) > u'_1(D, D) > u'_1(C, D)
$$
και
$$
u'_2(C, D) > u'_2(C, C) > u'_2(D, D) > u'_2(D, C)
$$

Ξεκινόντας για τον παίκτη 1, έχουμε:
$$
\begin{cases}
  u'_1(D, C) > u'_1(C, C) \\
  u'_1(C, C) > u'_1(D, D) \\
  u'_1(D, D) > u'_1(C, D)
\end{cases}
$$
Δηλαδή:
$$
\begin{cases}
  3 > 2 + 2\alpha \\
  2 + 2\alpha > 1 + \alpha \\
  1 + \alpha > 3\alpha
\end{cases}
$$
Από αυτό παίρνουμε
$$
\begin{cases}
  \alpha < \frac{1}{2} \\
  \alpha > -\frac{1}{2} \\
  \alpha < \frac{1}{2}
\end{cases}
\,.$$
Αντίστοιχα για τον παίκτη 2 θα έχουμε συμμετρικά το ίδιο αποτέλεσμα. Άρα,
$\alpha \in [0, \frac{1}{2})$.

\subsection{Ερώτηση (ii)}

Για να είναι το προφίλ $(C, C)$ σημείο ισορροπίας, θα πρέπει για κάθε παίκτη $i$
να ισχύει
$$
    u_i(\textbf{p}_i, \textbf{p}_{-i}) \geq u_i(\textbf{p}'_i, \textbf{p}_{-i}) \,.
$$

Συγκεκριμένα, για σε αυτό το παίγνιο θα πρέπει να ισχύει
\begin{align*}
    & u_1(C, C) \geq u_1(D, C) \land u_2(C, C) \geq u_2(C, D) \\
    \Leftrightarrow\ & 2 + 2a \geq 3 \land 2 + 2a \geq 3 \\
    \Leftrightarrow\ & a \geq \frac{1}{2} \,.
\end{align*}
